\SourceLeaf{167}
An arrangement that might rather seem admissible would be to set the fixed part of the pension for widows of the second class below 200 Rthl., for example at 150 Rthl., while allowing them an equal share in the supplementary dividend. If, however, one considers

that the loss of contributions begins immediately upon the formation of the second class, and then increases continually as more new extraordinary professors are appointed;

that among the 51 members of the widows' fund there are now 20 extraordinary professors, and therefore, if a similar proportion continues in future, a part of the increase in contributions assumed in the main plan, gradually growing to \(\tfrac{2}{5}\), will fail to materialize; whereas, on the other hand,

compensation from the lower pensions that come into payment will in all probability begin only after a long time and, given the very low mortality at the age at which most extraordinary professors usually are, will also occur only rarely on average,

it remains very doubtful whether setting the fixed part of the pension at 150 Rthl. for the second class would suffice even to compensate the fund in full for the loss through the lower contributions. Still less could such a measure be assumed to be decidedly to the fund's \textit{advantage}. [One would be entitled to do so only if either the contributions in the second class were kept as high as in the first, or the pension were reduced considerably below 150 Rthl.; but either course would seem so harsh that the commission cannot endorse it.]

Another measure which might sometimes bring the fund some advantage would be to repeal the provision in Article 10 of the regulations by which a widow's pension lapses entirely on her remarriage. The purpose of this provision can only have been to secure a gain for the fund on the assumption that such cases would occur fairly often. But this purpose is almost entirely missed, since experience shows remarriages of widows in these circumstances to be almost unheard of. (It has already been noted in the memorandum that only one such case has occurred in more than 100 years.) The arrangement in other widows' funds, by which the widow's pension is merely suspended during a second marriage but revives if the remarried woman becomes a widow for the second time, is undoubtedly more suitable. Compare the statutes of the Hof- und Civil-Diener Witwenkasse, §. 25, the Prediger-Witwenkasse, §. 24, and the Schullehrer-Witwenkasse, §. 22. An exception would probably have to be made where the second husband is again a professor here, since it seems inadmissible to grant one widow (or one family) two pensions from our widows' fund. On the other hand, it would seem advisable to dispense with further restrictions on revival, and in particular to restore the full pension for a professor's widow even where the widow, through her second marriage, for example to a clergyman or another state servant, received a pension from a \textit{different} public fund. For it must be considered that a change to the existing state of affairs can be expected to have an effect only to the extent that it is not itself cancelled by exceptions. It is self-evident, moreover, that where there are children under 20 from the first marriage, the orphan's pension would have to continue during the second marriage to the same extent as before.

\Needspace{5\baselineskip}
\pdfbookmark[2]{Note on the commission report.}{witwen-iii-note}
\begin{center}\textit{Note on the commission report.}\end{center}
\noindent\hspace*{\parindent}When the expressions balance and deficit are used in connection with a state's financial budget, the former means the comparison of expenditure and income expected for one year, or for a small number of years forming a financial period, according to a provisional
\SourceLeaf{168}
estimate; deficit means the difference when expenditure amounts to a larger sum than income. In such a context, surpluses and deficit are direct opposites, and each excludes the other of itself.

The calculation made for the widows' fund in the second division of the memorandum differs from such a calculation in two essential respects.

1) It compares expenditure and income for neither one year nor a few years, but includes both for the whole unlimited future, as far and as accurately as it is possible to estimate it in advance.

2) It sums not the sums of money themselves, at the amounts in which they will be spent or received (which would in any case be meaningless because of the absence of a limit), but their values discounted or reduced to the initial epoch at a specified rate of interest. In this sense even a series of expenditure or income continuing without limit can nevertheless be summed to a definite finite result. This possibility is readily apparent if one considers how enormously a sum of money dwindles when discounted over a very long period: for example, a payment due after 422 years in the amount of 7000 Rthl. has, at an interest rate of \(3\tfrac{1}{2}\) per cent, a present value of only \textit{one pfennig}. The same result is reached still more easily by considering that the present monetary value of a payment to be made every year in the same amount without end is nothing other than the capital sum which, at the chosen rate, yields an equal amount of interest annually.

It is already self-evident that if, on the factual basis of such a calculation, expenditure will be greater in the future than now, while income other than interest on capital may not be assumed in the calculation to increase, the calculation will show a deficit if income does not now yield a surplus over expenditure. But one cannot simply assert the converse, that if income now yields surpluses over expenditure, there will be \textit{no} deficit in the total calculation: this requires not merely the existence of surpluses, but surpluses of sufficient size. If a correct total balance calculation has shown a deficit as its final result, this establishes that the present annual surpluses are too small to cover the shortfalls to come later. The correct calculation has already taken the present temporary surpluses into account as parts of the balance, and they must not be credited to the institution \textit{twice}.

Just as false as the belief that present temporary surpluses will reduce a presently existing deficit is the idea that this deficit will have been extinguished as soon as the assets have increased by an amount equal to it. Such a conclusion is admissible only in the single case where the increase in assets is brought about \textit{immediately}, and by \textit{external} inflows not already included in the balance. Suppose, for example, that after 20 years the assets of the widows' fund had increased (perhaps partly through new inflows, but not of a recurring kind) by 17520 Rthl.: even then the deficit will not thereby have been removed, but may possibly be greater than now. For the balance on 1 October 1865 will result from comparing the fund's assets then existing with the amount of all expenditure to come from then on, discounted to \textit{that} epoch. This amount will be much greater than the corresponding amount for 1 October 1845, for the simple reason that in 1865 the time of the much greater annual expenditure will have drawn so much closer; indeed, very probably, such expenditure will already have begun to a considerable extent.

